\begin{proof}

\begin{enumerate}[label={$\langle1\rangle$\arabic*.},leftmargin=*]
\item Use the comparison maps of P4.52, and let $D_1$ be the
effective total conductor relative to $\omega_{S_1}$.
Write $n=\operatorname{ord}_H(E)>0$ and
$a=\operatorname{coeff}_H(D_1)$. The pole identity is
$2=n+a$, so $(n,a)=(2,0)$ or $(1,1)$.

\item Suppose $H$ dominates a reduced component $L$ of the
projective cubic tangent curve. At its generic point
a transverse Weierstrass presentation has degree
$m$, the multiplicity of that component in the cubic.
Factoring over the completed valuation field of the
exceptional coordinate $r$ gives the exact degree
identity
\[
 \sum_{H'}\operatorname{ord}_{H'}(r)
                 [\mathbb C(H'):\mathbb C(L)]=m.             \tag{4.53a}
\]
The factors group the roots by valuation branches;
the two factors in each summand are the ramification
index and residue-field degree. Residue extensions
are separable in characteristic zero.

\item A positive-genus component of a cubic is a smooth
elliptic cubic of multiplicity one. The strict
transform is smooth along it, with $n=1,a=0$,
contrary to $2=n+a$. Every other reduced component
is rational. The finite map from $H$ then has degree
$h\ge2$. Equation (4.53a) excludes $n=2$, since
$m\le3$. Thus $n=a=1$.

\item At a generic point of $L$, the $h$ points on $H$
give distinct branches of the transverse plane germ,
each with conductor exponent one. The branch formula
\[
             c_i=2\delta_i+\sum_{j\ne i}I_{ij}
\]
forces $h=2$ and exactly two smooth branches meeting
once, with no additional branch. Each has $r$-order
one, so their total Weierstrass degree is two. Hence
$m=2$, and the tangent polynomial is $y^2\ell$,
where $y=0$ is the double line and $\ell$ a different
linear form.

\item Choose complementary coordinates $x,z$. The
restriction $\ell(x,z)$ is nonzero; any $y$ term in
$\ell$ has higher weight for the weights used below.
In a generic blowup chart along the double line,
write $v$ for the transverse coordinate. The
strict-transform equation begins
\[
                 \ell v^2+r(f_4|_L)+\cdots.
\]
A nonzero $f_4|_L$ would give ramification order two
and residue degree one. Therefore $f_4$ is divisible
by $y$. The two branches meet once, so after
$v=ru$ their distinct leading slopes solve
\[
                         \ell u^2+G_3u+G_5=0,
\]
where $G_3,G_5$ are forms of degrees three and five
in $x,z$. The discriminant
$\Delta=G_3^2-4\ell G_5$ is nonzero and nonsquare
in $\mathbb C(\mathbb P^1)$. The actual residue-field
extension of $H\to L$ is
$\mathbb C(\mathbb P^1)(\sqrt\Delta)$.

\item Consequently the original equation has weighted
leading polynomial
\[
          P_5=\ell(x,z)y^2+G_3(x,z)y+G_5(x,z),               \tag{4.53b}
\]
of degree five for weights $(1,2,1)$ on $(x,y,z)$,
with no smaller-weight terms. This identifies the
normalized projective component birationally with
$H$ through its function field.

\item Let $D=\{P_5=0\}\subset\mathbb P(1,2,1)$.
The quadratic is irreducible over $\mathbb C(x/z)$,
since the degree-two residue cover $H$ is connected.
The only possible further factor is a common
vertical factor of $\ell,G_3,G_5$. Since $\ell$
is linear, either $P_5$ is irreducible or
$P_5=\ell Q_4$, with $Q_4$ monic quadratic in $y$.
Both alternatives are reduced.

\item In the irreducible case, let $k_i$ be the
multiplicities of the distinct zeros of the
homogeneous degree-six discriminant. The normalized
double cover has genus
$(\#\{i:k_i\text{ odd}\}-2)/2$. Positive genus
therefore requires at least four odd multiplicities,
and gives
\[
              \sum_i(k_i-1)=6-\#\{i\}\le2.                  \tag{4.53c}
\]
Away from $\ell=0$ and the weighted vertex,
completing the quadratic identifies the projective
curve germ with $w^2-s^{k_i}$, up to a unit and
holomorphic coordinates. Its transverse Milnor
number is $k_i-1$. At $\ell=0$, if $G_3\ne0$ the
finite root is simple and the other point is at
the weighted vertex. If $G_3=0$, absence of a
common factor gives $G_5\ne0$, and $\Delta$ has
a simple zero. There is no additional singularity.
In the chart $y=1$ the vertex is the quotient of
$\mathbb C^2$ by $(x,z)\mapsto(-x,-z)$; the
lifted equation $\ell+G_3+G_5$ has nonzero
differential. The original polynomial also has
nonzero derivative in the $\ell$ direction on
the $y$-axis, so that axis is not critical.
The total transverse Milnor-number sum is thus
$M\le2$.

\item In the reducible case write
$G_3=\ell A_2$, $G_5=\ell A_4$. Then
$Q_4=y^2+A_2y+A_4$ has discriminant of degree four.
A connected positive-genus double cover requires
four distinct simple roots. Its normalization is
elliptic and $Q_4$ is smooth on the weighted plane;
it avoids the weighted vertex. The rational
component $\ell=0$ meets it in total length two,
read directly from the restricted monic quadratic.
Either there are two transverse intersections,
with total Milnor number two, or one tangency of
order two, with Milnor number three. There are no
other critical rays. Hence in all cases
\[
                              M\le3.                       \tag{4.53d}
\]

\item The weighted plane has Euler number three: its
strata $x\ne0$, $x=0,z\ne0$, and $x=z=0$ are
$\mathbb C^2$, $\mathbb C$, and a point.
In the irreducible case $e(D)\le e(H)\le0$.
In the reducible case
$e(D)=0+2-\#(Q_4\cap\{\ell=0\})\le1$.
Therefore
\[
               \chi\bigl(\mathbb P(1,2,1)\setminus D\bigr)
                         \ge2.                            \tag{4.53e}
\]

\item The weighted scalar action of a primitive fifth
root, with weights $(1,2,1)$, acts freely on
$F=\{P_5=1\}$. Its quotient is
$\mathbb P(1,2,1)\setminus D$: the only possible
nontrivial stabilizer of the weighted
$\mathbb C^*$ action is at the weighted vertex,
which belongs to $D$. This is an unramified
fivefold cover. The compact-core and character
calculation of P4.52 gives
\[
 \operatorname{rank}(T-I)
 \ge4\chi\bigl(\mathbb P(1,2,1)\setminus D\bigr)
 \ge8.
\]
The variation identity (4.52a) applies to every
one-dimensional hypersurface critical locus. Its
boundary mapping-torus estimate would give
$\operatorname{rank}(T-I)\le M\le3$ if the
intersection form were zero.

\item It is zero here by the actual anisotropic scaling
comparison. Set
\[
 D_\epsilon(x,y,z)=(\epsilon x,\epsilon^2y,\epsilon z).
\]
Then $\epsilon^{-5}f(D_\epsilon(x,y,z))$ converges
with derivatives to $P_5$. The compact-cycle flow
of P4.52 transports its pairings into the actual
local Milnor fibre, whose pairing vanishes by
P4.51. This contradiction excludes every dominating
$H$.

\item Thus $H$ is exceptional over the new normalization.
Its coefficient $a$ is positive, so $n=a=1$.
Adjunction gives
\[
                    0=2g(H)-2+(D_1-H)\cdot H.
\]
It follows that $H$ is elliptic and meets no other
positive pole component. A zero-coefficient
exceptional curve cannot meet a positive one:
its own intersection equation has nonnegative
canonical term on a minimal resolution and
nonnegative remaining terms, with a strictly
positive term at such a meeting. The connected
exceptional fibre is therefore $H$ alone and is
disjoint from the strict conductor. The
normalization is an isomorphism at its image;
another analytic hypersurface branch would
intersect this branch in a curve and create
conductor meeting $H$. Hence the entire
strict-transform germ is normal, with the single
smooth elliptic exceptional curve $H$. This is
precisely an isolated simple elliptic singularity.

\item The discriminant multiplicity partitions and the
exact rank bounds are recomputed in
\texttt{code/part4/cubic\_weighted.py} and its
executed output.

\end{enumerate}
\end{proof}
